\section*{الفصل \ref{ch:g9:inside-the-atom} --- داخل الذرة}

\begin{solution}{exo:g9:inside-the-atom:1}
في \omterm{def:g9:inside-the-atom:nucleus}{النواة}: \omterm{def:g9:inside-the-atom:proton}{البروتونات} (موجبة) \omterm{def:g9:inside-the-atom:proton}{والنيوترونات} (بلا شحنة). وحولها:
\omterm{def:g9:inside-the-atom:nucleus}{الإلكترونات} (سالبة).
\end{solution}

\begin{solution}{exo:g9:inside-the-atom:2}
\ce{^{16}_{8}O}: 8 \omterm{def:g9:inside-the-atom:proton}{بروتونات}، و8 \omterm{def:g9:inside-the-atom:proton}{نيوترونات}، و8 \omterm{def:g9:inside-the-atom:nucleus}{إلكترونات}.
\ce{^{27}_{13}Al}: 13، 14، 13. \ce{^{40}_{20}Ca}: 20، 20، 20.
\end{solution}

\begin{solution}{exo:g9:inside-the-atom:3}
فيها من \omterm{def:g9:inside-the-atom:nucleus}{الإلكترونات} (شحنة كل منها $-e$) بقدر ما فيها من \omterm{def:g9:inside-the-atom:proton}{البروتونات} (شحنة
كل منها $+e$): فتتلاشى الشحنات.
\end{solution}

\begin{solution}{exo:g9:inside-the-atom:4}
$A = 26 + 30 = 56$: \ce{^{56}_{26}Fe}.
\end{solution}

\begin{solution}{exo:g9:inside-the-atom:5}
نعم: لكليهما $Z = 17$، فكلاهما كلور. ولهما عددان كتليان مختلفان (18 و20
نيوترونًا): فهما نظيران.
\end{solution}

\begin{solution}{exo:g9:inside-the-atom:6}
لا: لهما عددان ذريان مختلفان (6 و7)، فهما عنصران مختلفان، الكربون
والنيتروجين. \omterm{def:g9:inside-the-atom:isotope}{والنظائر} لها $Z$ نفسه.
\end{solution}

\begin{solution}{exo:g9:inside-the-atom:7}
$56 \times 1.67 \times 10^{-27} \approx \qty{9.35e-26}{kg}$.
\end{solution}

\begin{solution}{exo:g9:inside-the-atom:8}
\omterm{def:g7:atoms-and-molecules:atom}{الذرة} في معظمها فراغ: \omterm{def:g9:inside-the-atom:nucleus}{والنواة} بالغة الصغر، فمعظم الجسيمات لا تلقى \omterm{def:g9:inside-the-atom:nucleus}{نواة} في
طريقها فتمر مباشرة.
\end{solution}

\begin{solution}{exo:g9:inside-the-atom:9}
$10^{-10} / 10^{-15} = 10^{5}$: أي أصغر بمئة ألف مرة.
\end{solution}

\begin{solution}{exo:g9:inside-the-atom:10}
الكيمياء تتوقف على \omterm{def:g9:inside-the-atom:nucleus}{الإلكترونات}، ولنظائر العنصر الواحد العدد نفسه من
\omterm{def:g9:inside-the-atom:nucleus}{الإلكترونات} (العدد $Z$ نفسه).
\end{solution}

\begin{solution}{exo:g9:inside-the-atom:11}
$26 \times 9.11 \times 10^{-31} = \qty{2.37e-29}{kg}$. والكسر:
$2.37 \times 10^{-29} / 9.35 \times 10^{-26} \approx 2.5 \times
10^{-4}$، أي نحو \qty{0.025}{\percent} من الكتلة.
\end{solution}

\begin{solution}{exo:g9:inside-the-atom:12}
6915 \omterm{def:g7:atoms-and-molecules:atom}{ذرة} من النحاس 63 و3085 \omterm{def:g7:atoms-and-molecules:atom}{ذرة} من النحاس 65. والمتوسط:
$(6915 \times 63 + 3085 \times 65)/10\,000 = 63.6$ نوكليونًا.
\end{solution}

\begin{solution}{pb:g9:inside-the-atom:1}
\textbf{1.} \ce{^{35}_{17}Cl} و\ce{^{37}_{17}Cl}.

\textbf{2.} الكلور 35: 17 بروتونًا، و18 نيوترونًا، و17 إلكترونًا.
الكلور 37: 17 بروتونًا، و20 نيوترونًا، و17 إلكترونًا.

\textbf{3.} لكليهما 17 بروتونًا، فكلاهما العنصر كلور؛ ولا يختلفان إلا في عدد
\omterm{def:g9:inside-the-atom:proton}{النيوترونات}، فهما نظيران.

\textbf{4.} نعم: فالكيمياء تتوقف على \omterm{def:g9:inside-the-atom:nucleus}{الإلكترونات}، وفي كليهما 17 إلكترونًا.

\textbf{5.} $35 \times 1.67 \times 10^{-27} = \qty{5.845e-26}{kg}$.

\textbf{6.} $17 \times 9.11 \times 10^{-31} = \qty{1.55e-29}{kg}$:
أي أقل من كتلة \omterm{def:g9:inside-the-atom:nucleus}{النواة} بنحو 4000 مرة، فهي مهملة.

\textbf{7.} $37 \times 1.67 \times 10^{-27} = \qty{6.179e-26}{kg}$.

\textbf{8.} $\qty{1}{g} = \qty{e-3}{kg}$؛ و$10^{-3} / 5.845 \times
10^{-26} \approx 1.71 \times 10^{22}$ \omterm{def:g7:atoms-and-molecules:atom}{ذرة}.

\textbf{9.} \qty{75.8}{\percent} و\qty{24.2}{\percent}.

\textbf{10.} $758 \times 5.845 \times 10^{-26} = \qty{4.430e-23}{kg}$؛
و$242 \times 6.179 \times 10^{-26} = \qty{1.495e-23}{kg}$.

\textbf{11.} $4.430 \times 10^{-23} + 1.495 \times 10^{-23} =
\qty{5.925e-23}{kg}$.

\textbf{12.} $5.925 \times 10^{-23} / 1000 \approx
\qty{5.93e-26}{kg}$: وهو متوسط كتلة \omterm{def:g7:atoms-and-molecules:atom}{ذرة} الكلور.
\end{solution}
